\section{Discussion}

Our results confirmed that fairly simple neural-network agents can
learn to coordinate in a referential game in which they need to
communicate about a large number of real pictures. 
They also
suggest that the meanings agents come to assign to symbols in this setup
 capture general conceptual properties of the objects depicted in the
image, rather than low-level visual properties. We also showed a path to grounding the communication in natural language by mixing the game with a supervised task.

In future work, encouraged by our preliminary experiments
with object naming, we want to study how to ensure that the emergent
communication stays close to human natural language. %We think that
Predictive learning should be retained as an
important building block of intelligent agents, focusing on teaching
them structural properties of language (e.g., lexical choice, syntax
 or style).  However, it is also important to learn the function-driven facets of language, 
such
as how to hold a conversation, and interactive games are a potentially fruitful method to achieve this goal. %How to best combine the two approaches in an
%integrated whole remains a central question for future research.

%In future work, we will consider more complex referential games, going
%beyond cheap talk to foster the development of compositional
%symbols. At the same time, encouraged by our preliminary experiments
%with object naming, we want to study how to insure that the emergent
%communication stays close to human natural language. We think that
%predictive learning from canned corpora should be retained as an
%important building block of intelligent agents, focusing on teaching
%them structural properties of language, such as lexical choice, syntax
%or style. In parallel, agents should learn function-driven behaviours, such
%as how to hold a conversation, in interactive games like the one we
%presented here. How to best combine the two approaches in an
%integrated whole remains a central question for future research.

% We believe that the most promising direction towards developing agents that can work together with us is to place emphasis on the interactive aspect of
% language.  After all, human-machine co-ordination, just like the human-human
% one,  is impossible without communication, and the bulk of communication happens through natural language.
% Along these lines, we advocate a research program that creates
% interactive environments between a community of artificial agents. The agents will be given tasks to solve (i.e., games similar to the referential game presented here), that will be designed in order to foster
% emergence of key properties of language. Language will not be the goal of these
% games, but rather the by-product of achieving situated goals in this
% environment. 

% Connecting the emerged communication to human natural language is a key challenge for this framework; supervised learning and interactive learning should both come together to achieve this goal. We think that predictive learning should be retained as an important building block of intelligent agents, but rather should 
% focused on teaching them primitive properties of language rather than complex behaviours, such as
% how to hold a conversation~\cite{Vinyals:Le:2015}.



